% ============================================================
% StatLab Beamer Theme — Custom Header
% Mirrors the color scheme and layout of custom.scss (RevealJS)
% ============================================================

% ── COLORS (matches custom.scss variables) ───────────────────
\usepackage{xcolor}
\usepackage{pifont}

\definecolor{statprimary}{HTML}{2c3e50}    % dark gray-blue (body text)
\definecolor{statsecondary}{HTML}{2980b9}  % blue (structure, links)
\definecolor{statyale}{HTML}{00356b}       % Yale blue (footer)

% Box backgrounds
% Two-color system: Yale blue + statsecondary blue dominate.
% Orange (#FF6600) reserved for key warnings / key findings only.
\definecolor{stataccbg}{HTML}{e8f4f8}      % light cyan bg  — default block body
\definecolor{statorange}{HTML}{FF6600}     % orange         — alert / key finding header
\definecolor{statorangebg}{HTML}{FFF2E6}   % pale orange bg — alert body
\definecolor{statexbg}{HTML}{e6eef5}       % pale Yale-blue bg — example / Sam block body
\definecolor{statproof}{HTML}{95a5a6}      % gray           — proof (rarely used)
\definecolor{statproofbg}{HTML}{f9f9f9}    % light gray bg

% ── FONT: CMU Serif (LaTeX default) ─────────────────────────
% Switch Beamer from its default sans-serif to the classic serif.
\usefonttheme{serif}
% For XeLaTeX, explicitly select the OpenType CMU Serif face:
\setmainfont{CMU Serif}

% ── BEAMER COLORS ────────────────────────────────────────────

% Frame chrome
\setbeamercolor{frametitle}{bg=statprimary, fg=white}
\setbeamercolor{title}{fg=statprimary}
\setbeamercolor{subtitle}{fg=statsecondary}
\setbeamercolor{author}{fg=statprimary!70}
\setbeamercolor{date}{fg=statprimary!70}
\setbeamercolor{normal text}{fg=statprimary, bg=white}
\setbeamercolor{structure}{fg=statsecondary}

% Bullets and lists
\setbeamercolor{item}{fg=statsecondary}
\setbeamercolor{subitem}{fg=statsecondary!70}
\setbeamercolor{enumerate item}{fg=statsecondary}
\setbeamercolor{enumerate subitem}{fg=statsecondary!70}

% Bullet shape: filled circle (level 1), open circle (level 2), dash (level 3)
\setbeamertemplate{itemize item}{\small$\bullet$}
\setbeamertemplate{itemize subitem}{\small$\circ$}
\setbeamertemplate{itemize subsubitem}{\small\textendash}

% Default block — blue, mirrors .theorem / .lemma / .corollary
\setbeamercolor{block title}{fg=white, bg=statsecondary}
\setbeamercolor{block body}{bg=stataccbg, fg=statprimary}

% Alert block — orange (#FF6600), reserved for key warnings / findings
\setbeamercolor{block title alerted}{fg=white, bg=statorange}
\setbeamercolor{block body alerted}{bg=statorangebg, fg=statprimary}

% Example block — Yale blue, used for Sam callbacks and contextual examples
\setbeamercolor{block title example}{fg=white, bg=statyale}
\setbeamercolor{block body example}{bg=statexbg, fg=statprimary}

% ── FONT SIZES ───────────────────────────────────────────────
\setbeamerfont{title}{size=\LARGE, series=\bfseries}
\setbeamerfont{subtitle}{size=\large}
\setbeamerfont{frametitle}{size=\large, series=\bfseries}
\setbeamerfont{block title}{size=\normalsize, series=\bfseries}
\setbeamerfont{footline}{size=\scriptsize}

% ── FRAME TITLE ──────────────────────────────────────────────
\setbeamertemplate{navigation symbols}{}

\setbeamertemplate{frametitle}{%
  \nointerlineskip%
  \begin{beamercolorbox}[%
    wd=\paperwidth, ht=2.6ex, dp=0.9ex, leftskip=0.9em%
  ]{frametitle}%
    \usebeamerfont{frametitle}\insertframetitle%
  \end{beamercolorbox}%
}

% ── HEADLINE: presentation title, top-right (hidden on title slide) ──
\setbeamercolor{headline}{fg=statprimary!55, bg=white}
\setbeamertemplate{headline}{%
  \ifnum\insertframenumber>1%
    \begin{beamercolorbox}[wd=\paperwidth, ht=2.2ex, dp=0.5ex]{headline}%
      \hfill%
      \usebeamercolor[fg]{headline}{\fontsize{6.5}{8}\selectfont\inserttitle}%
      \hspace*{0.9em}%
    \end{beamercolorbox}%
  \else%
    \vspace{2.7ex}% reserve same height on title slide so frame title aligns
  \fi%
}

% ── FOOTER: branding + slide number (no logo) ────────────────
\setbeamertemplate{footline}{%
  {\color{statprimary!25}\hrule height 0.4pt}%
  \begin{beamercolorbox}[wd=\paperwidth, ht=2.4ex, dp=1.4ex]{footline}%
    \hspace*{0.9em}%
    \textcolor{statyale}{\fontfamily{qpl}\selectfont Yale Library $|$ StatLab}%
    \hfill%
    \textcolor{statprimary!50}{\insertframenumber\,/\,\inserttotalframenumber}%
    \hspace*{0.9em}%
  \end{beamercolorbox}%
}

% ── IMAGE CONSTRAINTS ────────────────────────────────────────
% Keep images within the content area so they never bleed into the footer.
% max-width = full text column; max-height = 65 % of the frame text height.
\usepackage{graphicx}
\setkeys{Gin}{width=\linewidth, height=0.65\textheight, keepaspectratio}

% ── TITLE PAGE ───────────────────────────────────────────────
\setbeamertemplate{title page}{%
  \vspace*{2.5em}%
  \begin{beamercolorbox}[sep=0pt, left]{title}%
    \usebeamerfont{title}\usebeamercolor[fg]{title}\inserttitle\par%
  \end{beamercolorbox}%
  \vspace*{0.5em}%
  \begin{beamercolorbox}[sep=0pt, left]{subtitle}%
    \usebeamerfont{subtitle}\usebeamercolor[fg]{subtitle}\insertsubtitle\par%
  \end{beamercolorbox}%
  \vspace*{1.8em}%
  \begin{beamercolorbox}[sep=0pt, left]{author}%
    \usebeamerfont{author}\usebeamercolor[fg]{author}\insertauthor\par%
  \end{beamercolorbox}%
  \vspace*{0.4em}%
  \begin{beamercolorbox}[sep=0pt, left]{date}%
    \usebeamerfont{date}\usebeamercolor[fg]{date}\insertdate\par%
  \end{beamercolorbox}%
}

% ── CUSTOM ENVIRONMENTS ──────────────────────────────────────
% Use these inside raw LaTeX blocks (```{=latex}) in your .qmd.
% Each mirrors the matching CSS class from the RevealJS theme.

% Definition — Yale blue (consistent with two-color palette)
\newenvironment{definitionblock}[1][Definition]{%
  \begingroup%
  \setbeamercolor{block title}{fg=white, bg=statyale}%
  \setbeamercolor{block body}{bg=statexbg, fg=statprimary}%
  \begin{block}{#1}%
}{%
  \end{block}%
  \endgroup%
}

% Lemma — blue, mirrors .lemma (same as default block)
\newenvironment{lemmablock}[1][Lemma]{%
  \begingroup%
  \begin{block}{#1}%
}{%
  \end{block}%
  \endgroup%
}

% Corollary — blue, mirrors .corollary (same as default block)
\newenvironment{corollaryblock}[1][Corollary]{%
  \begingroup%
  \begin{block}{#1}%
}{%
  \end{block}%
  \endgroup%
}

% Proof with QED square — gray, mirrors .proof
\newenvironment{proofblock}[1][Proof]{%
  \begingroup%
  \setbeamercolor{block title}{fg=white, bg=statproof}%
  \setbeamercolor{block body}{bg=statproofbg, fg=statprimary}%
  \begin{block}{\textit{#1}}%
  \itshape%
}{%
  \hfill$\square$%
  \end{block}%
  \endgroup%
}

% ── Workshop slide packages ───────────────────────────────────────────────
\usepackage{listings}
\usepackage{forest}
\usepackage{tikz}
\usetikzlibrary{arrows,arrows.meta,positioning,shapes.geometric,calc}
\usepackage{booktabs}
\usepackage{multicol}
\usepackage{array}
\usepackage{colortbl}   % \cellcolor / \rowcolor for highlighting table lookups

% ── Workshop color definitions ────────────────────────────────────────────
% Inline semantic colors (used in tables/code for before/after comparisons)
\definecolor{errorred}{RGB}{180,30,30}       % red  — bad examples, error states
\definecolor{highlightgreen}{RGB}{0,130,80}  % green — good examples, cached targets
\definecolor{commentgray}{RGB}{120,120,120}  % gray — code comments
\definecolor{warnorange}{RGB}{200,100,0}     % warm orange — string literals in code
\definecolor{lightblue}{RGB}{220,235,255}    % light blue — tikz fill

% ── Listings: tiny font, StatLab color accents ────────────────────────────
\lstset{
  basicstyle=\ttfamily\tiny,
  breaklines=true,
  frame=single,
  framerule=0.4pt,
  framesep=3pt,
  xleftmargin=8pt,
  xrightmargin=4pt,
  showstringspaces=false,
  commentstyle=\color{commentgray}\itshape,
  keywordstyle=\color{statsecondary}\bfseries,
  stringstyle=\color{warnorange},
  numbers=left,
  numberstyle=\tiny\color{gray!60},
  numbersep=5pt,
  tabsize=2,
  aboveskip=3pt,
  belowskip=3pt,
}

% ── Forest dirtree style ──────────────────────────────────────────────────
\forestset{
  dirtree/.style={
    for tree={
      font=\ttfamily\scriptsize,
      grow'=0,
      child anchor=west,
      parent anchor=south,
      anchor=west,
      calign=first,
      inner sep=1pt,
      outer sep=0pt,
      edge path={
        \noexpand\path [draw, \forestoption{edge}]
        (!u.south west) +(7.5pt,0) |- (.child anchor)
        \forestoption{edge label};
      },
      before typesetting nodes={
        if n=1{insert before={[,phantom]}}{}
      },
      fit=band,
      before computing xy={l=15pt},
    },
  }
}

% ── Quarto code block font size override ──────────────────────────────────
% Load fancyvrb explicitly so \DefineVerbatimEnvironment is always available,
% whether or not pandoc includes code blocks that trigger its auto-load.
\usepackage{fancyvrb}
\AtBeginDocument{%
  % Source code (syntax-highlighted input)
  \DefineVerbatimEnvironment{Highlighting}{Verbatim}{%
    commandchars=\\\{\},%
    fontsize=\small%
  }%
  % Printed results (plain verbatim) — matched to the input so a chunk
  % and its output read at one consistent size.
  \RecustomVerbatimEnvironment{verbatim}{Verbatim}{fontsize=\small}%
}

% ── Two-column shorthand colors for code comparisons ─────────────────────
\newcommand{\badlabel}{{\color{errorred}\textbf{$\boldsymbol{\times}$ Before}}}
\newcommand{\goodlabel}{{\color{highlightgreen}\textbf{\ding{51}~After}}}

% ── ORR Workshop rebuild macros (2025) ────────────────────────────────────

% Yale blue alias (for \color{yale})
\colorlet{yale}{statyale}

% Figure placeholder — reserves whitespace and flags human action needed
\newcommand{\figplaceholder}[1]{%
  \begin{center}%
  \fbox{\begin{minipage}{0.82\textwidth}%
  \centering\vspace{1.4cm}%
  {\large\color{gray} [FIGURE: #1]}\\[0.3cm]%
  {\small\color{gray}\itshape See [HUMAN INSTRUCTION] comment in source}%
  \vspace{1.4cm}%
  \end{minipage}}%
  \end{center}%
}

% Key finding callout (red alert block)
\newcommand{\findingbox}[1]{%
  \begin{alertblock}{Key Finding}%
  #1%
  \end{alertblock}%
}

% Title-less callout — example-block coloring, no title bar.
% Use where a block would otherwise invent a second heading under the frame title.
\newenvironment{notebox}{%
  \begin{center}%
  \begin{minipage}{0.95\textwidth}%
  \begin{beamercolorbox}[wd=\textwidth, sep=6pt]{block body example}%
}{%
  \end{beamercolorbox}%
  \end{minipage}%
  \end{center}%
}

% Inline pain/solution markers
\newcommand{\pain}[1]{{\color{errorred}\textbf{Problem:}} #1}
\newcommand{\solutiontag}[1]{{\color{yale}\textbf{Solution:}} #1}

% Parker Lab scenario callback — green example block
\newcommand{\samcallback}[1]{%
  \begin{exampleblock}{Parker Lab --- Sam's Problem}%
  \small #1%
  \end{exampleblock}%
}

% ── Named lstlisting styles ────────────────────────────────────────────────
% Override the global \tiny with \scriptsize for readability on slides.

\lstdefinestyle{Rcode}{%
  language=R,%
  basicstyle=\ttfamily\scriptsize,%
  backgroundcolor=\color{gray!8},%
  frame=single,%
  breaklines=true,%
  breakatwhitespace=false,%
  showstringspaces=false,%
  numbers=none,%
}

\lstdefinestyle{bashcode}{%
  language=bash,%
  basicstyle=\ttfamily\scriptsize,%
  backgroundcolor=\color{gray!8},%
  frame=single,%
  breaklines=true,%
  numbers=none,%
}

\lstdefinestyle{Pythoncode}{%
  language=Python,%
  basicstyle=\ttfamily\scriptsize,%
  backgroundcolor=\color{gray!8},%
  frame=single,%
  breaklines=true,%
  numbers=none,%
}

\lstdefinestyle{jsoncode}{%
  language={},%
  basicstyle=\ttfamily\scriptsize,%
  backgroundcolor=\color{gray!8},%
  frame=single,%
  breaklines=true,%
  numbers=none,%
  morestring=[b]",%
  stringstyle=\color{warnorange},%
}

\lstdefinestyle{yamlcode}{%
  language={},%
  basicstyle=\ttfamily\scriptsize,%
  backgroundcolor=\color{gray!8},%
  frame=single,%
  breaklines=true,%
  numbers=none,%
  morecomment=[l]{\#},%
  commentstyle=\color{commentgray}\itshape,%
}

\lstdefinestyle{textcode}{%
  language={},%
  basicstyle=\ttfamily\scriptsize,%
  backgroundcolor=\color{gray!8},%
  frame=none,%
  breaklines=true,%
  numbers=none,%
}
